% ================================================
% 文件: 04_线材.tex
% 章节: 第5章 线材
% 所属部分: 电子元器件与材料
% 版本: v1.0
% ================================================
% 本章目录结构（树状）:
% \chapter{线材}
% ├── \section{线材基础}
% ├── \section{导线类型与特性}
% ├── \section{电缆与传输线}
% ├── \section{连接器与接头}
% └── \section{线材选型与应用}
% ================================================

\chapter{线材}
\label{chap:wiring_materials}

\section{线材基础}
\label{sec:wiring_fundamentals}

线材是电子设备中用于传输电能和信号的导体材料。线材的主要参数包括：

\begin{itemize}
    \item \textbf{导体材料}：铜、铝、银等，铜是最常用的导体材料
    \item \textbf{截面积}：常用AWG（美国线规）或mm²表示
    \item \textbf{绝缘材料}：PVC、PE、PTFE等
    \item \textbf{额定电压}：导线安全工作的最高电压
\end{itemize}

\section{导线类型与特性}
\label{sec:wire_types}

常见的导线类型包括：

\begin{itemize}
    \item \textbf{单芯线}：单根导体，用于固定安装
    \item \textbf{多股线}：多根细导体绞合，柔韧性好
    \item \textbf{漆包线}：表面涂漆绝缘，用于绕制线圈
    \item \textbf{电磁线}：用于电机、变压器绕组
\end{itemize}

导线的直流电阻：$R = \rho\frac{l}{S}$，其中$\rho$为电阻率，$l$为长度，$S$为截面积。

\section{电缆与传输线}
\label{sec:cables_transmission_lines}

电缆是由多根导线组成的传输线路，常见类型：

\begin{itemize}
    \item \textbf{同轴电缆}：内导体、绝缘层、外导体、护套，用于高频信号传输
    \item \textbf{双绞线}：两根导线绞合，用于差分信号传输
    \item \textbf{屏蔽电缆}：带有屏蔽层，抗干扰能力强
    \item \textbf{电力电缆}：用于传输大功率电能
\end{itemize}

同轴电缆的特性阻抗通常为50Ω或75Ω。

\section{连接器与接头}
\label{sec:connectors}

连接器用于连接导线或设备，常见类型：

\begin{itemize}
    \item \textbf{RF连接器}：SMA、BNC、N型等，用于射频信号
    \item \textbf{音频连接器}：XLR、RCA、3.5mm耳机插头
    \item \textbf{电源连接器}：插头插座、端子排
    \item \textbf{电路板连接器}：排针排母、板对板连接器
\end{itemize}

连接器的关键参数包括接触电阻、绝缘电阻、插拔次数等。

\section{线材选型与应用}
\label{sec:wiring_selection}

选型时需考虑：

\begin{itemize}
    \item \textbf{电流容量}：根据载流量选择线径
    \item \textbf{频率特性}：高频应用需考虑趋肤效应
    \item \textbf{环境条件}：温度、湿度、化学腐蚀等
    \item \textbf{机械强度}：抗拉强度、柔韧性要求
\end{itemize}

高频应用中，趋肤效应使电流集中在导体表面，需使用多股绞合线或镀银导线。